\section{Tritic orientation, compatible prolongation and double projection}
\label{sec:trit-projection}

The generated character of the real coordinates is expressed in the relation
between arithmetic depth and evaluation. The primitive datum is an admissible
history of operations with its residues, quotients and orientations. An
Archimedean reading associates increasingly narrow intervals with its prefixes;
an incidence reading preserves boundary relations of the same state.
The double projection of the corpus \cite{hmtI,HMT-IX} brings these two operations together
without identifying a scalar evaluation with the entire history. The
article uses this antecedent to study the subsequent realizations
of the dodecaphasic record.

\subsection{Positional cylinders and determination of the real coordinate}

For a ternary block \(b=(b_1,\ldots,b_6)\), with representatives
\(b_j\in\{0,1,2\}\) of \(\mathbb F_3\), the integer evaluation
\(\nu(b)=\sum_{j=1}^6b_j3^{6-j}\) belongs to
\(\{0,\ldots,728\}\). If \(b_{n+1}\) is the next admissible block,
positional prolongation satisfies
\[
 N_{n+1}=729N_n+\nu(b_{n+1}),\qquad
 a_n=\frac{N_n}{729^n},\qquad b_n^+=\frac{N_n+1}{729^n}.
\]
Here \(N_0\) is the integer part determined by the regional reader. The symbol
\(b_n^+\) denotes a real endpoint and is distinct from the ternary block \(b\).
The representation cylinder is \([a_n,b_n^+)\); its closure is used
in the argument establishing existence of the limit. Carry preserves the passage from
a prefix to its predecessor through Euclidean division.

\begin{proposition}[Limit of cylinder centres]
For a compatible prolongation, the cylinder closures are
nested and determine a unique coordinate \(x\). If
\(c_n=(a_n+b_n^+)/2\), then
\[
 x=\lim_{n\to\infty}a_n=\lim_{n\to\infty}b_n^+
 =\lim_{n\to\infty}c_n,\qquad
 |x-c_n|\leq\frac1{2\,729^n}.
\]
\end{proposition}
\begin{proof}
The inequality \(0\leq\nu(b_{n+1})<729\) places the subcylinder in
the closure of its predecessor. Its endpoints are rational, its width is
\(729^{-n}\), and it remains inside the first bounded closed interval.
Nesting and decreasing width produce a unique Cauchy
class of endpoints. In the Archimedean realization, this class has the
representative \(x\); the distance to the midpoint is at most half
the width. The boundary convention assigns a representation when two
positional expansions represent the same limit.
\end{proof}

The arithmetic mean thus enters a precise evaluation: the
centre of each interval bounding the coordinate. The limit preserves the
exactness of the value as the width vanishes. The enriched
representation additionally preserves the prefix sequence and its incidence.
The numerical projection and the exceptional projection are two maps
with different codomains; one preserves the limiting value and the other transports
the relations enabling the lattice construction. The passage to
\(\RR\) is thus situated in the realization of the reader, after the
arithmetic generation of the data it evaluates.

To make compatibility explicit, let \(p_{n+1,n}\) be the restrictions
of the enriched state, with \(x_n=p_{n+1,n}(x_{n+1})\). If \(I_n(x_n)\)
is the numerical cylinder and \(\mathcal B_n(x_n)\) the boundary incidence,
the two transports simultaneously preserve
\[
 \overline{I_{n+1}(x_{n+1})}\subseteq\overline{I_n(x_n)},\qquad
 r_{n+1,n}\mathcal B_{n+1}(x_{n+1})=\mathcal B_n(x_n).
\]
The second equality uses the incidence maps of the corpus;
Proposition~\ref{prop:lf-incidence-naturality} makes their realization
through transported frames and codes explicit. Its five
consubstantial constructions---spectral, solenoidal, gauge, cohomological
and determinantal---belong to that same compatible system. The subsequent
geometric study preserves this unity of origin by specifying which
of its observations each theorem uses.

\subsection{Oriented trajectories and the temporality of update}

Let \(x\xrightarrow{\rm adm}y\) be the TPK admissibility relation on
the enriched space \(\widehat X\). A trajectory parametrized by
a discretely ordered set \(\mathcal T\) is
\[
 \gamma:\mathcal T\longrightarrow\widehat X,
 \qquad\gamma(t)\xrightarrow{\rm adm}\gamma(t^+).
\]
The compatibility relation determines the transitions; the parametrization
allows them to be ordered. In the space \(E=\mathcal T\times\widehat X\),
with \(p(t,x)=t\), the trajectory determines the section
\(\widetilde\gamma(t)=(t,\gamma(t))\), which satisfies
\(p\circ\widetilde\gamma=\id_{\mathcal T}\).
This distinction preserves the difference between a state, its trajectory and
the coordinate used to describe its evolution.

The tritic signature \(\tau\in\{+1,0,-1\}\) types the generators through
\[
 J_\tau^2=-\tau I.
\]
In the temporal reader of the corpus, \(+1\) represents return and memory;
\(0\), threshold update; and \(-1\), the opening of conjugate
prolongations. Their interpretation as past, present and future refers
to the functions of the preserved antecedent, the effective transition and the
admissible continuation. The three tenses thus acquire an
operator-theoretic content on histories.

If \(\mathcal P_n(x)\) is the set of admissible continuations of
length \(n\), deleting the final step defines
\(r_n:\mathcal P_{n+1}(x)\to\mathcal P_n(x)\). An unbounded continuation
is a collection \((\gamma_n)\) with
\(r_n(\gamma_{n+1})=\gamma_n\). Restriction reconstructs antecedents;
the update \(U_t\) acts at the transition; prolongation maintains
the prefix boundary. The existence of these collections is used in
the surviving domains fixed by the prolongation theorem of the
corpus.

On a finite signature \(\sigma=(\tau_1,\ldots,\tau_n)\), oriented
inversion is
\[
 C_{\rm rev}\sigma=(-\tau_n,\ldots,-\tau_1),
 \qquad C_{\rm rev}^2\sigma=\sigma.
\]
The conjugate trajectory also preserves the sheet,
position and boundary transformations of its transport. Bidirectionality is therefore
expressed in the full history, while this formula gives its signature.
The phase return of \(\Gamma_9\) admits a memory increment:
states with identical phase may occupy different depths.

\subsection{Quadratic algebras and activity of the parabolic threshold}

In a fixed regime, conjugation of \(aI+bJ_\tau\) changes the sign
of \(b\), and its algebraic norm is
\[
 (aI+bJ_\tau)(aI-bJ_\tau)=(a^2+\tau b^2)I.
\]
The elliptic, parabolic and hyperbolic types correspond to the three
quadratic relations. Their local groups are represented by
\[
 e^{tJ_{+1}}=\cos t\,I+\sin t\,J_{+1},\quad
 e^{tJ_0}=I+tJ_0,\quad
 e^{tJ_{-1}}=\cosh t\,I+\sinh t\,J_{-1}.
\]
The proof separates even and odd terms in the exponential series; in the
nilpotent type, only degrees zero and one remain. Each collection
satisfies the composition law in its parameter and the inversion
\(e^{-tJ_\tau}=(e^{tJ_\tau})^{-1}\).

In particular, the zero tritic value permits a nonzero generator
with square zero. For
\(N=\left(\begin{smallmatrix}0&1\\0&0\end{smallmatrix}\right)\),
\[
 e^{tN}(q,p)=(q+tp,p).
\]
The threshold has an effective linear action on its fibre. Its parameter,
its displacement and the memory of the transition remain determined
even though the regime signature is zero. This is the precise content
allowing the present to be developed narratively as an active update.

The positive elliptic form has circular level sets in orthonormal
coordinates and elliptical ones after a linear transformation; the hyperbolic form
has two isotropic directions and hyperbolic level sets; the parabolic type
has a degenerate form and unipotent transport. Spherical geometry
additionally requires the homogeneous space or metric realizing it. The
preceding classification fixes the generators that may participate in that
realization and distinguishes algebraic type from the curvature of a
subsequent space.

The mode calendar makes the connection with self-scaling concrete. The rule
\(+1:m\mapsto\Sigma\), \(-1:m\mapsto\Pi\), \(0:m\mapsto m\)
applied to the cyclic block \((+1,-1,0)\) produces
\((\Sigma,\Pi,\Pi)\). Counting its three transitions in the order
\((\Sigma,\Pi)\) determines
\begin{equation}
 F_{\rm av}=\begin{pmatrix}0&1\\1&1\end{pmatrix}.
 \label{eq:geo-fav}
\end{equation}
Its positive eigenray, normalized as \((1,x)^\mathsf T\), satisfies
\(x=\lambda\), \(1+x=\lambda x\), and hence
\(x^2-x-1=0\). The positive root is the self-scaling coordinate obtained
as \(\varphi\). The matrix arises from the calendar of operations; the
radical expression and the following exponential identities are its
subsequent realizations.

% Articulación posterior a la generación de pi, phi y e.
% Contenido acumulativo; no reemplaza el desarrollo previo del TRIT.
\subsection{Exponential realizations of the generated coordinates}
\label{geo-euler-trit-articulacion}

The already determined regional coordinates allow the concrete
realizations of the quadratic TRIT collection to be recognized. The order is
substantive: construction of the regions precedes evaluation of an
exponential at the parameters \(\pi\) or \(2\log\varphi\). These expressions
identify the role of the coordinates in operators constructed from
APP and TPK; they do not participate in selecting the prefixes that produce them.

\subsubsection{Ternary cycle, nilpotent transport and sheet incidence}

Multiplication \(a(r)=4r\) in \(\mathbb Z/9\mathbb Z\) has order
\(3\). On the units it determines the orbits \(\{1,4,7\}\) and
\(\{2,8,5\}\). In the augmentation module of either orbit, formed by
integer combinations whose coefficients sum to zero, the basis
\((e_r-e_{4r},e_{4r}-e_{7r})\) gives
\[
 C=\begin{pmatrix}0&-1\\1&-1\end{pmatrix},
 \qquad C^2+C+I=0.
\]
The restricted form has Gram matrix
\(\bigl(\begin{smallmatrix}2&-1\\-1&2\end{smallmatrix}\bigr)\).
In residue evaluations modulo \(9\), the action distinguishes the sheets:
\(S(ax,ay)=aS(x,y)\), whereas \(P(ax,ay)=a^{-1}P(x,y)\).
The cyclic realization thus preserves a contragredient orientation
between sum and product. The quotients of the integer lift are preserved
in the corresponding arithmetic record.

Elementary parabolic transport is represented by
\[
 N=\begin{pmatrix}0&1\\0&0\end{pmatrix}.
\]
The areal--volumetric incidence already constructed in \eqref{eq:geo-fav} gives
\[
 F_{\rm av}=\begin{pmatrix}0&1\\1&1\end{pmatrix},
 \qquad A=F_{\rm av}^2=\begin{pmatrix}1&1\\1&2\end{pmatrix}.
\]
The three operators act on declared real spaces: the augmentation
module tensored with \(\mathbb R\), the plane of nilpotent transport and
the modal incidence plane. Having the same matrix dimension does not identify
these spaces or their coordinates.

\begin{proposition}[Concrete generators of the three types]
\label{geo-euler-trit-generadores}
The operators
\[
 J_{\mathrm e}=\frac{2C+I}{\sqrt3},\qquad
 J_{\mathrm p}=N,\qquad
 J_{\mathrm h}=\frac{2A-3I}{\sqrt5}
\]
satisfy, respectively,
\[
 J_{\mathrm e}^{\,2}=-I,\qquad
 J_{\mathrm p}^{\,2}=0,\qquad
 J_{\mathrm h}^{\,2}=I.
\]
\end{proposition}
\begin{proof}
The relation for \(C\) implies \((2C+I)^2=-3I\).
Direct multiplication gives \(N^2=0\).
Finally, \(A\) has trace \(3\) and determinant \(1\), so
\(A^2-3A+I=0\) and \((2A-3I)^2=5I\).
\end{proof}

\begin{proposition}[Three evaluations of the exponential identity]
\label{geo-euler-trit-evaluaciones}
In the preceding realizations,
\begin{equation}
 \boxed{
 \exp(\pi J_{\mathrm e})+I=0,\qquad
 \exp(J_{\mathrm p})=I+J_{\mathrm p},\qquad
 \exp(2\log\varphi\,J_{\mathrm h})=F_{\rm av}^{\,2}.}
 \label{geo-euler-trit-tres-evaluaciones}
\end{equation}
\end{proposition}
\begin{proof}
The elliptic specialization of the tritic exponential gives
\(\cos\pi\,I+\sin\pi\,J_{\mathrm e}=-I\).
Nilpotence removes exactly the terms of degree at least \(2\)
in the second expression. For the third, \(\varphi^2=\varphi+1\) implies
\[
 \cosh(2\log\varphi)=\frac32,\qquad
 \sinh(2\log\varphi)=\frac{\sqrt5}{2}.
\]
The exponential therefore equals
\(\frac32I+\frac{\sqrt5}{2}J_{\mathrm h}=A\).
\end{proof}

The first equality represents the elliptic half-turn; the full turn
satisfies \(\exp(2\pi J_{\mathrm e})=I\). The second retains a nonzero
linear term: the parabolic regime expresses transport, not absence of
evolution. The third realizes a unimodular hyperbolic step. Its
eigenvalues are \(\varphi^2\) and \(\varphi^{-2}\); one direction expands and
the other contracts while preserving the determinant. The three evaluations
use different parameters of a common exponential collection.

In that collection, \(e\) enters through the scalar evaluation
\(\exp(I)=eI\) and the law \(\exp(sI)\exp(tI)=\exp((s+t)I)\).
Its role is not restricted to the parabolic type. The coordinate \(e\), previously
obtained from its region, is recognized here as the unit-parameter evaluation
of the exponential; this identity does not replace its coinductive genealogy.

\subsubsection{Polynomial resolution of the regimes}

The three realizations admit a compact presentation on the direct
sum of their spaces. Define
\[
 \mathbb J=J_{\mathrm e}\oplus J_{\mathrm p}\oplus J_{\mathrm h},
 \qquad Q=\mathbb J^2.
\]
The preceding relations imply
\[
 \mathbb J^6=\mathbb J^2,\qquad Q^3=Q.
\]
Using \(Q\) distinguishes this operator square from the
dodecaphasic record \(K\).

\begin{proposition}[Polynomial regime projectors]
\label{geo-euler-trit-proyectores}
The maps
\[
 p_{\mathrm e}=\frac{Q^2-Q}{2},\qquad
 p_{\mathrm p}=I-Q^2,\qquad
 p_{\mathrm h}=\frac{Q^2+Q}{2}
\]
are mutually orthogonal idempotents, sum to \(I\) and recover the three
summands of the representation.
\end{proposition}
\begin{proof}
The three polynomials take the values \((1,0,0)\), \((0,1,0)\) and
\((0,0,1)\) at the points \(-1,0,+1\), respectively. The differences
between their squares and themselves, as well as their cross-products, are
multiples of \(X^3-X\). Substituting \(Q\), which annihilates that polynomial, proves
the identities. On the three blocks, \(Q\) acts as \(-I,0,I\).
\end{proof}

It follows that
\begin{align}
 \exp(t\mathbb J)
 ={}&p_{\mathrm e}(\cos t\,I+\sin t\,\mathbb J)
   +p_{\mathrm p}(I+t\mathbb J)\notag\\
   &+p_{\mathrm h}(\cosh t\,I+\sinh t\,\mathbb J).
 \label{geo-euler-trit-exponencial-total}
\end{align}
Each summand preserves its quadratic relation, and conjugation reverses
its evolution. The full representation brings together the three regimes, but does not
by itself define a transition operator between them. Such a transition
requires the TPK transport maps with their domains, orientation and
memory. The exponential articulation allows their types to be recognized without
replacing that dynamics by an identification of the fibres.

% PROCEDENCIA FOCAL (no impresa)
% Resultado recuperado: tres generadores y tres evaluaciones exponenciales.
% Propietario completo:
% BASE/manuscrito/sections/hmt/09b_algebras_cuadraticas.tex:289-516.
% Antecedente causal de C y acción contragrediente:
% BASE/manuscrito/sections/hmt/03d_esqueleto_ciclotomico_app.tex:13-47,145-246.
% Los tres tipos y la orientación están en:
% BASE/manuscrito/sections/hmt/02_trit_geometrias.tex:83-196.
% Resolución polinómica recuperada:
% output/INVESTIGACION_HOLONOMIA_NONADICA_CRISTAL_APERIODICO_20260903/
% ALGEBRA_HOLONOMICA_UNIVERSAL_Y_EULER_TRITICO.md, secciones 8-11 y 22.
% Los cuatro propietarios se leyeron completos. BASE designa:
% /Users/ruben/Documents/HMT2/
% EL_CIERRE_HOLOGRAFICO_DEL_INFINITO_SUCESOR_102_CAPITULOS_2026-08-28_EN_TRABAJO/
% No se incorpora la identidad P9(alpha)=delta del documento especializado:
% esa identificación no pertenece a este bloque.
% Tampoco se identifica una dirección nilpotente con la vacancia electrónica,
% ni se presume una representación global del transporte por la suma directa.
% COMPROBACIONES: Python estándar, enteros y Fraction.
% C^2+C+I=0; (2C+I)^2=-3I; N^2=0; (2Fav^2-3I)^2=5I.
% 81 pares de residuos verifican las dos acciones contragredientes.
% Los 3 proyectores se verificaron en Q=-1,0,+1; prueba general en texto.
% 6 labels únicos y entornos equilibrados. Sin compilación.

